% ======================================================================
% Topic 3 -- Counting and Probability -- part c (problems 31--40)
% ======================================================================
\clearpage
\dividerbar{Extension and challenge}

% ---------------------------------------------------------------- Problem 31
\begin{probbox}{31}{\relax}
A pencil case holds $6$ red pens, $9$ blue pens and $5$ black pens, identical except for the colour of the ink. One pen is taken at random.
\begin{enumerate}[label=\alph*),leftmargin=1.5em,itemsep=0pt,topsep=2pt]
\item Find the probability of taking a blue pen.
\item Find the probability of taking a pen that is \emph{not} black.
\item How many more black pens must be put into the case so that the probability of taking a black pen is $\dfrac12$?
\end{enumerate}
\end{probbox}
\probsol{
\textbf{Understanding the problem.} Each pen is one equally likely outcome. Parts a and b are direct probability calculations; part c is a ``backwards'' problem: when black pens are added, \emph{both} the number of black pens \emph{and} the total number of pens increase.
\step{1}{Total number of pens}
$6+9+5=20$ pens.
\step{2}{Part a}
\begin{recallbox}$P=\dfrac{\text{number of favourable outcomes}}{\text{number of all outcomes}}$.\end{recallbox}
$P(\text{blue})=\dfrac{9}{20}$.
\step{3}{Part b -- complement}
\begin{recallbox}$P(\text{not }A)=1-P(A)$.\end{recallbox}
$P(\text{black})=\dfrac{5}{20}=\dfrac14$, so $P(\text{not black})=1-\dfrac14=\dfrac34$.

\emph{Check by counting directly:} the pens that are not black are $6+9=15$ pens, $\dfrac{15}{20}=\dfrac34$ \checkmark.
\step{4}{Part c -- choose an unknown and form an equation}
Let $x$ be the number of black pens to be added ($x$ is a natural number). Only black pens are added, so afterwards:
\begin{itemize}[leftmargin=1.2em,itemsep=0pt,topsep=1pt]
\item the number of black pens is $5+x$;
\item the total number of pens is $20+x$.
\end{itemize}
The condition in the problem: $P(\text{black})=\dfrac{5+x}{20+x}=\dfrac12$.
\step{5}{Solve the equation}
\begin{recallbox}For $b\ne0$, $d\ne0$: $\dfrac ab=\dfrac cd$ exactly when $a\times d=b\times c$ (property of proportions).\end{recallbox}
From $\dfrac{5+x}{20+x}=\dfrac12$ we get $2\times(5+x)=1\times(20+x)$, that is
\[10+2x=20+x\ \Rightarrow\ 2x-x=20-10\ \Rightarrow\ x=10.\]
So $10$ black pens must be added.

\emph{Remark:} $P(\text{black})=\dfrac12$ also means that the number of black pens \emph{equals} the number of non-black pens. There are always $15$ non-black pens, so there must be $15$ black pens, i.e.\ we add $15-5=10$ -- in agreement with the solution above.
\begin{warnbox}Many students think ``to go from $\dfrac14$ to $\dfrac12$, double the black pens, so add $5$''. But when $5$ black pens are added the total also rises to $25$, and the probability is $\dfrac{10}{25}=\dfrac25$, not yet $\dfrac12$.\end{warnbox}
\verify{Adding $10$ black pens: $15$ black pens out of $20+10=30$ pens; $P(\text{black})=\dfrac{15}{30}=\dfrac12$~\checkmark.}
\ans{a) $\dfrac{9}{20}$;\quad b) $\dfrac34$;\quad c) add $10$ black pens.}
}

% ---------------------------------------------------------------- Problem 32
\begin{probbox}{32}{\relax}
\begin{minipage}[t]{0.6\linewidth}\vspace{0pt}
A spinner is divided into four sectors with central angles $180^\circ$ (red), $90^\circ$ (blue), $60^\circ$ (yellow) and $30^\circ$ (purple), as shown. The spinner is spun once.
\begin{enumerate}[label=\alph*),leftmargin=1.5em,itemsep=0pt,topsep=2pt]
\item Find the probability that the arrow lands on blue.
\item Find the probability that the arrow lands on yellow or purple.
\item Leo says: ``There are $4$ colours, so each colour has probability $\dfrac14$.'' Is Leo right or wrong? Why?
\end{enumerate}
\end{minipage}\hfill
\begin{minipage}[t]{0.36\linewidth}\vspace{0pt}
\centering\begin{tikzpicture}[scale=1.05]
\fill[AnsRed!45] (0,0)--(90:1.2) arc (90:270:1.2)--cycle;
\fill[tmBlue!45] (0,0)--(270:1.2) arc (270:360:1.2)--cycle;
\fill[StarGold!70] (0,0)--(0:1.2) arc (0:60:1.2)--cycle;
\fill[tmPurple!45] (0,0)--(60:1.2) arc (60:90:1.2)--cycle;
\draw[Ink,line width=1pt] (0,0) circle (1.2);
\foreach \a in {90,270,0,60}{\draw[Ink,line width=0.8pt] (0,0)--(\a:1.2);}
\node[font=\footnotesize] at (180:0.65){$180^\circ$};
\node[font=\footnotesize] at (315:0.7){$90^\circ$};
\node[font=\footnotesize] at (30:0.75){$60^\circ$};
\node[font=\scriptsize] at (75:0.9){$30^\circ$};
\draw[-{Stealth[length=6pt]},line width=1.4pt,Ink] (0,0)--(130:0.9);\fill[Ink] (0,0) circle (1.5pt);
\end{tikzpicture}
\end{minipage}
\end{probbox}
\probsol{
\textbf{Understanding the problem.} The sectors have \emph{different} sizes, so the four sectors are \emph{not} equally likely. We cannot use ``$1$ sector out of $4$''. Instead, the chance of the arrow stopping in a sector is proportional to the size of the sector, that is, to its central angle.
\begin{recallbox}For a spinner with unequal sectors: $P(\text{sector})=\dfrac{\text{central angle of the sector}}{360^\circ}$. The whole circle has $360^\circ$.\end{recallbox}
\step{1}{Check the data}
$180^\circ+90^\circ+60^\circ+30^\circ=360^\circ$ -- exactly one full circle \checkmark.
\step{2}{Part a}
$P(\text{blue})=\dfrac{90}{360}=\dfrac14$.

\emph{Another way to think:} we can divide the spinner into $12$ equal small sectors of $30^\circ$ each. Then red takes $6$ small sectors, blue $3$, yellow $2$ and purple $1$; these $12$ small sectors are equally likely. $P(\text{blue})=\dfrac{3}{12}=\dfrac14$ \checkmark.
\step{3}{Part b}
The yellow and purple sectors do not overlap, so their total angle is $60^\circ+30^\circ=90^\circ$.
\[P(\text{yellow or purple})=\frac{90}{360}=\frac14\quad\Big(\text{or }\frac{2+1}{12}=\frac{3}{12}=\frac14\Big).\]
\step{4}{Part c}
Leo is \textbf{wrong}. The formula ``favourable divided by all'' only works when the outcomes are \emph{equally likely}. Here the four colours are not equally likely: the red sector is half of the spinner, so $P(\text{red})=\dfrac{180}{360}=\dfrac12$, while $P(\text{purple})=\dfrac{30}{360}=\dfrac{1}{12}$ -- very different.
\verify{$P(\text{red})+P(\text{blue})+P(\text{yellow})+P(\text{purple})=\dfrac{6}{12}+\dfrac{3}{12}+\dfrac{2}{12}+\dfrac{1}{12}=1$ \checkmark.}
\begin{parentbox}Cut a paper plate into $12$ equal ``slices of cake'' and colour them as in the problem. Counting the slices of each colour, your child will see at once why a big sector has a big probability.\end{parentbox}
\ans{a) $\dfrac14$;\quad b) $\dfrac14$;\quad c) Wrong, because the sectors are not equal (e.g.\ $P(\text{red})=\dfrac12$, $P(\text{purple})=\dfrac1{12}$).}
}

% ---------------------------------------------------------------- Problem 33
\begin{probbox}{33}{\relax}
The school canteen has this lunch menu: one main dish out of $3$ (chicken rice, spaghetti, pizza); one drink out of $2$ (apple juice, orange juice); one dessert out of $2$ (ice cream, yoghurt). Each lunch consists of one main dish, one drink and one dessert.
\begin{enumerate}[label=\alph*),leftmargin=1.5em,itemsep=0pt,topsep=2pt]
\item How many different lunches are there?
\item Mia is given a lunch at random. Find the probability that it includes spaghetti.
\item Find the probability that Mia's lunch includes both spaghetti and orange juice.
\end{enumerate}
\end{probbox}
\probsol{
\textbf{Understanding the problem.} Each lunch is made by three successive choices: main dish, drink, dessert. ``Given at random'' means all the different lunches are equally likely.
\step{1}{Part a -- multiplication rule}
\begin{recallbox}Multiplication rule: if a task has stages with $m$, $n$ and $p$ ways, there are $m\times n\times p$ ways altogether.\end{recallbox}
Number of lunches: $3\times2\times2=12$.
\begin{center}\begin{tikzpicture}[grow'=right,level distance=2.6cm,font=\sffamily\footnotesize,
  every node/.style={draw=tmBlue,fill=tmBlue!8,rounded corners=2pt,inner sep=2pt},
  level 1/.style={sibling distance=2.4cm},level 2/.style={sibling distance=1.2cm},level 3/.style={sibling distance=0.6cm},
  edge from parent/.style={draw=Ink!60}]
\node[fill=Ink!10,draw=Ink]{Lunch}
  child{node{chicken rice} child{node{apple juice} child{node{ice cream}} child{node{yoghurt}}} child{node{orange juice} child{node{ice cream}} child{node{yoghurt}}}}
  child{node[fill=tmOrange!25,draw=tmOrange]{spaghetti} child{node{apple juice} child{node{ice cream}} child{node{yoghurt}}} child{node[fill=tmOrange!25,draw=tmOrange]{orange juice} child{node[fill=tmOrange!25,draw=tmOrange]{ice cream}} child{node[fill=tmOrange!25,draw=tmOrange]{yoghurt}}}}
  child{node{pizza} child{node{apple juice} child{node{ice cream}} child{node{yoghurt}}} child{node{orange juice} child{node{ice cream}} child{node{yoghurt}}}};
\end{tikzpicture}\end{center}
The tree has $3$ big branches, each splitting into $2$ branches, each of which splits into $2$ again: $12$ end points, exactly $12$ lunches.
\step{2}{Part b}
Lunches with spaghetti: the main dish is fixed as spaghetti, drink $2$ ways, dessert $2$ ways: $1\times2\times2=4$ lunches (the middle branch of the tree).
$P(\text{spaghetti})=\dfrac{4}{12}=\dfrac13$.
\begin{tipbox}A quick way to think: the main dish is chosen at random from $3$ dishes, and each dish appears in the same number of lunches, so $P(\text{spaghetti})=\dfrac13$. The drink and dessert do not matter.\end{tipbox}
\step{3}{Part c}
Main dish spaghetti, drink orange juice, only the dessert is left to choose: $1\times1\times2=2$ lunches (orange in the tree).
$P(\text{spaghetti and orange juice})=\dfrac{2}{12}=\dfrac16$.
\verify{List them: (spaghetti, orange juice, ice cream) and (spaghetti, orange juice, yoghurt) -- exactly $2$ lunches \checkmark.}
\ans{a) $12$ lunches;\quad b) $\dfrac13$;\quad c) $\dfrac16$.}
}

% ---------------------------------------------------------------- Problem 34
\begin{probbox}{34}{\relax}
A hall has $20$ rows of seats. The first row has $12$ seats, and each row behind has $2$ seats more than the row just in front of it.
\begin{enumerate}[label=\alph*),leftmargin=1.5em,itemsep=0pt,topsep=2pt]
\item How many seats are in row $20$?
\item Which row has exactly $30$ seats?
\item How many seats does the hall have altogether?
\end{enumerate}
\end{probbox}
\probsol{
\textbf{Understanding the problem.} The numbers of seats in the rows are $12,14,16,\dots$ -- an \emph{evenly spaced sequence} with step $d=2$. We need the $20$th term, the position of the term $30$, and the sum of $20$ terms.
\step{1}{Part a -- find the pattern}
\begin{center}
\begin{tabular}{|c|c|c|c|c|c|}\hline
Row & $1$ & $2$ & $3$ & $4$ & $k$\\ \hline
Seats & $12$ & $12+2$ & $12+2\times2$ & $12+3\times2$ & $12+(k-1)\times2$\\ \hline
\end{tabular}
\end{center}
From row $1$ to row $k$ there are $k-1$ ``jumps'', each adding $2$ seats. Row $20$: $12+19\times2=12+38=50$ seats.
\begin{warnbox}Not $12+20\times2=52$: from row $1$ to row $20$ there are only $19$ increases (the ``off by one'' error).\end{warnbox}
\step{2}{Part b}
We need $12+(k-1)\times2=30$, so $(k-1)\times2=18$, $k-1=9$, $k=10$. Row $10$ has $30$ seats.
\step{3}{Part c -- sum of an evenly spaced sequence}
\begin{recallbox}Sum of an evenly spaced sequence $=\dfrac{(\text{first}+\text{last})\times\text{number of terms}}{2}$ (Gauss's pairing method).\end{recallbox}
Total number of seats: $\dfrac{(12+50)\times20}{2}=\dfrac{62\times20}{2}=620$ seats.
\begin{proofbox}Pair row $1$ with row $20$: $12+50=62$; row $2$ with row $19$: $14+48=62$; \ldots; row $10$ with row $11$: $30+32=62$. There are $10$ pairs of $62$ seats each, so the total is $10\times62=620$.\end{proofbox}
\verify{Another way: if every row had $12$ seats there would be $20\times12=240$ seats; the ``extra'' is $0+2+4+\dots+38=2\times(0+1+\dots+19)=2\times190=380$; total $240+380=620$ \checkmark.}
\ans{a) $50$ seats;\quad b) row $10$;\quad c) $620$ seats.}
}

% ---------------------------------------------------------------- Problem 35
\begin{probbox}{35}{\stars{1}}
A box contains only red and blue marbles. If one marble is drawn at random, the probability that it is red is $\dfrac35$. If $5$ more blue marbles are put into the box, the probability of drawing a red marble drops to $\dfrac12$. How many red marbles and how many blue marbles were in the box at the start?
\end{probbox}
\probsol{
\textbf{Understanding the problem.} We do not know the numbers of marbles, only two probabilities (two ratios). We choose unknowns for the numbers of red and blue marbles at the start, turn each condition into an equation, and solve the system of two equations. Note: adding blue marbles does \emph{not} change the number of red marbles, while the total increases by $5$.
\step{1}{Choose the unknowns}
Let $x$ be the number of red marbles and $y$ the number of blue marbles at the start ($x$, $y$ non-zero natural numbers). At the start the box holds $x+y$ marbles.
\step{2}{Turn the first condition into an equation}
$P(\text{red})=\dfrac{x}{x+y}=\dfrac35$. Cross-multiplying: $5x=3(x+y)$, i.e.\ $5x=3x+3y$, or
\[2x=3y.\qquad(1)\]
\begin{recallbox}$\dfrac ab=\dfrac cd$ (with $b$, $d$ non-zero) exactly when $a\times d=b\times c$.\end{recallbox}
\step{3}{Turn the second condition into an equation}
After adding $5$ blue marbles there are still $x$ red marbles, and the total is $x+y+5$. So $\dfrac{x}{x+y+5}=\dfrac12$. Cross-multiplying: $2x=x+y+5$, or
\[x=y+5.\qquad(2)\]
(That is: after adding, the number of red marbles \emph{equals} the number of blue marbles.)
\step{4}{Solve the system by substitution}
Substitute $x=y+5$ from (2) into (1):
\[2(y+5)=3y\ \Rightarrow\ 2y+10=3y\ \Rightarrow\ 3y-2y=10\ \Rightarrow\ y=10.\]
Then $x=y+5=10+5=15$.
\begin{tipbox}\emph{Using proportions:} from (1), $\dfrac x3=\dfrac y2$; put $\dfrac x3=\dfrac y2=k$, so $x=3k$, $y=2k$. Substituting into (2): $3k=2k+5$, so $k=5$, $x=15$, $y=10$.\end{tipbox}
\verify{At the start: $15$ red, $10$ blue, $P(\text{red})=\dfrac{15}{25}=\dfrac35$ \checkmark. Adding $5$ blue: $15$ red, $15$ blue, $P(\text{red})=\dfrac{15}{30}=\dfrac12$ \checkmark.}
\begin{warnbox}Do not try to ``subtract probabilities'': $\dfrac35-\dfrac12=\dfrac1{10}$ tells us nothing directly about the number of marbles. A probability is a \emph{ratio}; go back to actual quantities.\end{warnbox}
\ans{At the start there were $15$ red and $10$ blue marbles.}
}

% ---------------------------------------------------------------- Problem 36
\needspace{10\baselineskip}
\begin{probbox}{36}{\stars{1}}
Recall: a palindrome is a number that reads the same forwards and backwards (for example $3443$, $7007$).
\begin{enumerate}[label=\alph*),leftmargin=1.5em,itemsep=0pt,topsep=2pt]
\item How many four-digit palindromes are there?
\item How many of them are even?
\item A four-digit palindrome is chosen at random. Find the probability that it is even.
\item Show that every four-digit palindrome is divisible by $11$.
\end{enumerate}
\end{probbox}
\probsol{
\textbf{Understanding the problem.} A four-digit number $\overline{abcd}$ is a palindrome when reading it backwards, $\overline{dcba}$, gives the same number: $d=a$ and $c=b$. So a four-digit palindrome has the form $\overline{abba}$, and only \emph{two} digits are chosen freely: $a$ and $b$.
\step{1}{Part a}
\begin{recallbox}Multiplication rule; the first digit of a number cannot be $0$.\end{recallbox}
\begin{center}\begin{tabular}{|l|c|c|c|c|}\hline
Digit & first ($a$) & second ($b$) & third & fourth\\ \hline
Condition & from $1$ to $9$ & from $0$ to $9$ & equals $b$ & equals $a$\\ \hline
Number of choices & $9$ & $10$ & $1$ & $1$\\ \hline
\end{tabular}\end{center}
$a$: $9$ ways ($1$ to $9$); $b$: $10$ ways; the last two digits are forced to equal $b$ and $a$. There are $9\times10=90$ numbers.
\step{2}{Part b}
A number is even when its ones digit is even. The ones digit is $a$, so $a$ must be even \emph{and} non-zero: $a\in\{2,4,6,8\}$ -- $4$ ways. $b$ still has $10$ ways. There are $4\times10=40$ even palindromes.
\begin{warnbox}The digit $0$ is even, but $a=0$ is not allowed because $a$ is also the first digit. So there are $4$, not $5$, choices for $a$.\end{warnbox}
\step{3}{Part c}
$P(\text{even})=\dfrac{40}{90}=\dfrac49$.
\step{4}{Part d -- proof}
\begin{proofbox}Write $\overline{abba}$ using place values:
$\overline{abba}=1000a+100b+10b+a=1001a+110b$.
Since $1001=11\times91$ and $110=11\times10$,
$\overline{abba}=11\times91\times a+11\times10\times b=11\times(91a+10b)$,
which is divisible by $11$.\end{proofbox}
\verify{Examples: $3443=11\times313$; $7007=11\times637$; $1221=11\times111$ \checkmark. Counting check: the palindromes starting with $1$ are $1001,1111,\dots,1991$ -- exactly $10$ numbers; nine first digits give $90$ numbers \checkmark.}
\ans{a) $90$;\quad b) $40$;\quad c) $\dfrac49$;\quad d) $\overline{abba}=11\times(91a+10b)$.}
}

% ---------------------------------------------------------------- Problem 37
\needspace{10\baselineskip}
\begin{probbox}{37}{\stars{1}}
Using the four digits $0,1,2,3$, four-digit numbers with all digits different are made (each digit used exactly once).
\begin{enumerate}[label=\alph*),leftmargin=1.5em,itemsep=0pt,topsep=2pt]
\item How many such numbers can be made?
\item How many of them are even?
\item One of the numbers from part a) is chosen at random. Find the probability that it is even.
\end{enumerate}
\end{probbox}
\probsol{
\textbf{Understanding the problem.} Each number is an arrangement of the four digits $0,1,2,3$ in the four places (thousands, hundreds, tens, ones), with the condition that the thousands digit is not $0$.
\step{1}{Part a -- choose the restricted position first}
\begin{recallbox}Multiplication rule. A position with a condition is chosen first.\end{recallbox}
\begin{center}\begin{tabular}{|l|c|c|c|c|}\hline
Place & thousands & hundreds & tens & ones\\ \hline
Condition & not $0$ & $3$ digits left & $2$ digits left & last digit\\ \hline
Number of choices & $3$ & $3$ & $2$ & $1$\\ \hline
\end{tabular}\end{center}
Thousands: $3$ ways ($1$, $2$ or $3$). Hundreds: $3$ digits left (now $0$ may be used): $3$ ways. Tens: $2$ ways. Ones: $1$ way. There are $3\times3\times2\times1=18$ numbers.
\begin{warnbox}Forgetting the condition ``thousands digit not $0$'' gives $4\times3\times2\times1=24$ -- which includes $6$ ``numbers'' starting with $0$, such as $0123$, that are really three-digit numbers.\end{warnbox}
\step{2}{Part b -- split into cases by the ones digit}
An even number ends in $0$ or $2$. These two cases differ because the digit $0$ affects the thousands place.
\begin{itemize}[leftmargin=1.2em,itemsep=1pt,topsep=1pt]
\item \textbf{Case 1: ones digit $0$.} The three digits $1,2,3$ go freely into the other three places (no more worry about $0$ in front): $3\times2\times1=6$ numbers.
\item \textbf{Case 2: ones digit $2$.} $0,1,3$ remain. Thousands not $0$: $2$ ways ($1$ or $3$); hundreds: $2$ ways; tens: $1$ way. There are $2\times2\times1=4$ numbers.
\end{itemize}
\begin{recallbox}Addition rule: for separate cases, add the numbers of ways in each case.\end{recallbox}
Number of even numbers: $6+4=10$.
\step{3}{Part c}
$P(\text{even})=\dfrac{10}{18}=\dfrac59$.
\verify{List the even numbers. Case 1: $1230,1320,2130,2310,3120,3210$ ($6$ numbers). Case 2: $1032,1302,3012,3102$ ($4$ numbers). Total $10$ \checkmark. Odd numbers (ending in $1$ or $3$): each case is like Case 2, with $4$ numbers, total $8$; and $10+8=18$ \checkmark.}
\ans{a) $18$ numbers;\quad b) $10$ even numbers;\quad c) $\dfrac59$.}
}

% ---------------------------------------------------------------- Problem 38
\begin{probbox}{38}{\stars{1}}
\begin{minipage}[t]{0.62\linewidth}\vspace{0pt}
The figure shows a rectangular grid of $2$ rows and $3$ columns of squares.
\begin{enumerate}[label=\alph*),leftmargin=1.5em,itemsep=0pt,topsep=2pt]
\item How many rectangles (including squares) are there in the figure altogether?
\item How many squares are there?
\item One of these rectangles is chosen at random. Find the probability that it is a square.
\end{enumerate}
\end{minipage}\hfill
\begin{minipage}[t]{0.34\linewidth}\vspace{0pt}
\centering\begin{tikzpicture}[scale=0.75]
\foreach \x in {0,1,2,3}{\draw[Ink,line width=1pt] (\x,0)--(\x,2);}
\foreach \y in {0,1,2}{\draw[Ink,line width=1pt] (0,\y)--(3,\y);}
\end{tikzpicture}
\end{minipage}
\end{probbox}
\probsol{
\textbf{Understanding the problem.} This is a grid with full-length lines, so the ``choose two lines'' trick works. We will also count by size to check.
\step{1}{Part a -- choose two horizontal and two vertical lines}
\begin{recallbox}A rectangle in a grid is determined by $2$ horizontal and $2$ vertical lines. The number of ways to choose $2$ of $n$ lines is $\dfrac{n\times(n-1)}{2}$.\end{recallbox}
The grid has $3$ horizontal and $4$ vertical lines.
\begin{itemize}[leftmargin=1.2em,itemsep=0pt,topsep=1pt]
\item Choose $2$ horizontal lines: $\dfrac{3\times2}{2}=3$ ways.
\item Choose $2$ vertical lines: $\dfrac{4\times3}{2}=6$ ways.
\end{itemize}
By the multiplication rule: $3\times6=18$ rectangles.
\step{2}{Check with a table of sizes}
Write ``height $\times$ width'' (in squares). A shape $w$ squares wide has $4-w$ horizontal positions; $h$ squares high, $3-h$ vertical positions.
\begin{center}
\begin{tabular}{|c|c|c|c|c|c|c|c|}\hline
Size & $1\times1$ & $1\times2$ & $1\times3$ & $2\times1$ & $2\times2$ & $2\times3$ & Total\\ \hline
Number & $2\times3=6$ & $2\times2=4$ & $2\times1=2$ & $1\times3=3$ & $1\times2=2$ & $1\times1=1$ & $18$\\ \hline
\end{tabular}
\end{center}
Both methods give $18$ \checkmark.
\step{3}{Part b}
The squares are the $1\times1$ and $2\times2$ shapes: $6+2=8$ squares.
\begin{warnbox}There is no $3\times3$ square because the grid is only $2$ squares high. And do not forget the two $2\times2$ squares (left and right, overlapping in the middle column).\end{warnbox}
\step{4}{Part c}
$P(\text{square})=\dfrac{8}{18}=\dfrac49$.
\ans{a) $18$;\quad b) $8$;\quad c) $\dfrac49$.}
}

% ---------------------------------------------------------------- Problem 39
\begin{probbox}{39}{\stars{1}}
Four friends Anna, Ben, Chloe and David are seated at random on four chairs in a row.
\begin{enumerate}[label=\alph*),leftmargin=1.5em,itemsep=0pt,topsep=2pt]
\item How many seating arrangements are there?
\item Find the probability that Anna and Ben sit next to each other.
\item Find the probability that Anna and Ben sit at the two ends of the row.
\end{enumerate}
\end{probbox}
\probsol{
\textbf{Understanding the problem.} ``Seated at random'' means all arrangements are equally likely. We count all arrangements, then the favourable ones for each event. Number the chairs $1,2,3,4$ from left to right.
\step{1}{Part a}
\begin{recallbox}Seating $n$ people on $n$ chairs: $n\times(n-1)\times\dots\times2\times1$ ways (multiplication rule; each next chair has one person fewer to choose from).\end{recallbox}
Chair $1$: $4$ choices of person; chair $2$: $3$; chair $3$: $2$; chair $4$: $1$. There are $4\times3\times2\times1=24$ arrangements.
\step{2}{Part b -- treat Anna and Ben as one ``block''}
Glue Anna and Ben together into one ``block'' taking $2$ neighbouring chairs. The table shows that the block has only $3$ positions:
\begin{center}\begin{tabular}{|c|c|c|c|c|}\hline
 & Chair $1$ & Chair $2$ & Chair $3$ & Chair $4$\\ \hline
Position 1 & \multicolumn{2}{c|}{\cellcolor{tmOrange!20}Anna and Ben} & & \\ \hline
Position 2 & & \multicolumn{2}{c|}{\cellcolor{tmOrange!20}Anna and Ben} & \\ \hline
Position 3 & & & \multicolumn{2}{c|}{\cellcolor{tmOrange!20}Anna and Ben}\\ \hline
\end{tabular}\end{center}
\begin{itemize}[leftmargin=1.2em,itemsep=1pt,topsep=1pt]
\item Position of the block (Anna, Ben): chairs $1$--$2$, $2$--$3$ or $3$--$4$: $3$ ways.
\item Inside the block, Anna on the left or Ben on the left: $2$ ways.
\item Chloe and David on the $2$ remaining chairs: $2\times1=2$ ways.
\end{itemize}
Favourable arrangements: $3\times2\times2=12$. So $P(\text{next to each other})=\dfrac{12}{24}=\dfrac12$.
\begin{warnbox}Do not forget to multiply by $2$ for swapping Anna and Ben inside the block: ``Anna--Ben'' and ``Ben--Anna'' are two different arrangements.\end{warnbox}
\step{3}{Part c}
Anna and Ben on chairs $1$ and $4$: Anna on chair $1$ and Ben on chair $4$, or the other way round -- $2$ ways. Chloe and David on chairs $2$, $3$: $2$ ways. Favourable arrangements: $2\times2=4$.
$P(\text{both ends})=\dfrac{4}{24}=\dfrac16$.
\verify{Part b, counted another way: if Anna sits at an end (chair $1$ or $4$, $2$ possibilities) there is only $1$ chair next to Anna; if Anna sits in the middle (chair $2$ or $3$) there are $2$ chairs next to Anna. Number of (Anna, Ben) neighbouring chair pairs: $2\times1+2\times2=6$; times $2$ ways to seat Chloe and David: $12$ \checkmark.}
\begin{parentbox}With $4$ people, let your child use $4$ name cards and actually lay out all $24$ arrangements in order (starting with Anna, then starting with Ben\ldots). Listing once helps the child trust the multiplication rule.\end{parentbox}
\ans{a) $24$ arrangements;\quad b) $\dfrac12$;\quad c) $\dfrac16$.}
}

% ---------------------------------------------------------------- Problem 40
\begin{probbox}{40}{\stars{2}}
A box contains $4$ red marbles and $6$ blue marbles of the same size. Grace draws one marble at random, does \emph{not put it back}, and then draws a second marble at random.
\begin{enumerate}[label=\alph*),leftmargin=1.5em,itemsep=0pt,topsep=2pt]
\item Find the probability that both marbles are red.
\item Find the probability that the two marbles have the same colour.
\item Find the probability that the two marbles have different colours.
\end{enumerate}
\end{probbox}
\probsol{
\textbf{Understanding the problem.} Grace draws two marbles \emph{one after the other} without replacement. Imagine the marbles numbered $1$ to $10$ (marbles $1$--$4$ red, marbles $5$--$10$ blue). Each \emph{ordered pair} (first marble, second marble) is one outcome, and these outcomes are equally likely.
\step{1}{Count all outcomes}
\begin{recallbox}Multiplication rule: the first draw has $10$ ways; for each of them, the second draw has $9$ marbles left, so $9$ ways.\end{recallbox}
Number of outcomes: $10\times9=90$ ordered pairs.
\step{2}{Part a -- both red}
First draw red: $4$ ways. Then $3$ red marbles are left, second draw red: $3$ ways. Favourable outcomes: $4\times3=12$.
$P(\text{red, red})=\dfrac{12}{90}=\dfrac{2}{15}$.
\step{3}{Tree diagram with the numbers of ways}
\begin{center}\begin{tikzpicture}[grow'=right,level distance=3.1cm,font=\sffamily\footnotesize,
  every node/.style={draw=Ink!60,rounded corners=2pt,inner sep=2.5pt,fill=white},
  level 1/.style={sibling distance=2.2cm},level 2/.style={sibling distance=1.1cm},
  edge from parent/.style={draw=Ink!60}]
\node[fill=Ink!10]{Start}
  child{node[fill=AnsRed!15,draw=AnsRed]{Red: $4$ ways}
     child{node[fill=AnsRed!15,draw=AnsRed]{Red: $3$ ways $\Rightarrow 4\times3=12$}}
     child{node[fill=tmBlue!15,draw=tmBlue]{Blue: $6$ ways $\Rightarrow 4\times6=24$}}}
  child{node[fill=tmBlue!15,draw=tmBlue]{Blue: $6$ ways}
     child{node[fill=AnsRed!15,draw=AnsRed]{Red: $4$ ways $\Rightarrow 6\times4=24$}}
     child{node[fill=tmBlue!15,draw=tmBlue]{Blue: $5$ ways $\Rightarrow 6\times5=30$}}};
\end{tikzpicture}\end{center}
Check: $12+24+24+30=90$ \checkmark{} -- the four branches cover every outcome.
\step{4}{Part b -- same colour}
Same colour has two separate cases: (red, red) with $12$ outcomes and (blue, blue) with $30$ outcomes.
\begin{recallbox}Addition rule for cases that do not overlap.\end{recallbox}
Favourable outcomes: $12+30=42$. $P(\text{same colour})=\dfrac{42}{90}=\dfrac{7}{15}$.
\step{5}{Part c -- complement}
``Different colours'' is the complement of ``same colour'': $P(\text{different colours})=1-\dfrac{7}{15}=\dfrac{8}{15}$.
\verify{Direct count: (red, blue) $24$ and (blue, red) $24$, total $48$; $\dfrac{48}{90}=\dfrac{8}{15}$ \checkmark.}
\begin{warnbox}The most common mistake: acting as if all $10$ marbles are still there for the second draw, computing $\dfrac{4}{10}\times\dfrac{4}{10}$ or counting $4\times4$ red pairs. Because the marble is \emph{not put back}, the second draw has only $9$ marbles and only $3$ red ones.\end{warnbox}
\begin{tipbox}You can also count without order: the number of ways to choose $2$ of $10$ marbles is $\dfrac{10\times9}{2}=45$, the number of ways to choose $2$ of the $4$ red marbles is $\dfrac{4\times3}{2}=6$, so $P(\text{red, red})=\dfrac{6}{45}=\dfrac{2}{15}$ -- the same result. What matters is that the numerator and denominator are counted \emph{the same way} (both ordered or both unordered).\end{tipbox}
\ans{a) $\dfrac{2}{15}$;\quad b) $\dfrac{7}{15}$;\quad c) $\dfrac{8}{15}$.}
}
